\documentclass[11pt,a4paper]{article}

\usepackage[margin=1in]{geometry}
\usepackage{hyperref}
\usepackage{listings}
\usepackage{xcolor}
\usepackage{booktabs}
\usepackage{array}
\usepackage{parskip}
\usepackage{titlesec}
\usepackage{fancyhdr}
\usepackage{fontenc}
\usepackage{inputenc}

\hypersetup{
    colorlinks=true,
    linkcolor=blue,
    urlcolor=blue,
}

\lstdefinestyle{hexstyle}{
    basicstyle=\ttfamily\small,
    breaklines=true,
    frame=single,
    backgroundcolor=\color{gray!10},
    rulecolor=\color{gray!40},
    xleftmargin=1em,
    xrightmargin=1em,
}

\lstdefinestyle{gostyle}{
    language=C, % close enough for syntax highlighting
    basicstyle=\ttfamily\small,
    keywordstyle=\color{blue},
    commentstyle=\color{green!60!black},
    stringstyle=\color{orange},
    breaklines=true,
    frame=single,
    backgroundcolor=\color{gray!10},
    rulecolor=\color{gray!40},
    xleftmargin=1em,
    xrightmargin=1em,
}

\pagestyle{fancy}
\fancyhf{}
\rhead{Firmware Image Analysis}
\lhead{}
\rfoot{Page \thepage}

\title{Firmware Image Analysis}
\author{}
\date{March 2026}

\begin{document}

\maketitle
\tableofcontents
\newpage

% ==============================================================
\section{Overview}

This document describes the structure of the firmware images analyzed, the
approach taken to parse them, and how every byte region is accounted for.
The parser is written in Go and walks the image structurally, the same
way a bootloader would, rather than scanning for magic bytes.

% ==============================================================
\section{Tooling and AI Usage}

The parser was written over approximately five days of part-time work. AI
(Claude) was used in the following ways:

\begin{itemize}
    \item \textbf{Boilerplate generation} for repetitive struct parsing
        (e.g.\ the \texttt{binary.Read} sequences for inode types). The
        pattern was established manually; AI accelerated writing similar
        blocks.

    \item \textbf{Language gap filling}, the underlying systems
        programming concepts were already understood, but Go-specific
        idioms and simple go language functions required lookup at the start (first time using Go).

    \item \textbf{Struct scaffolding} for the \texttt{printFirmwareFile}
        function, to complete the menial work of printing out every attribute in a structure.
          
    \item \textbf{Write-up organization} for this document. The format
      	internals are well understood and I can answer questions about any
      	of it, but translating that into structured writing is a weaker
      	area. AI provided a starting framework that I then edited and
      	refined until satisfied.
    
    \item \textbf{Magic byte identification} during early reconnaissance.
      	I had not used binwalk before and initially preferred working with
      	tools I was already comfortable with. AI helped identify known
      	signatures directly in hex before I eventually adopted binwalk
      	as well.
    
	\item \textbf{Bouncing Ideas}, AI is useful for quickly sketching
      out a mockup to test a hypothesis, which significantly sped up
      working through theories about how the format is structured.

\end{itemize}

AI was \textbf{not} used to determine byte
offsets, understand structural metadata, or how the code in general should be written and organized. All format understanding came
from:

\begin{itemize}
    \item Manual hex inspection (HxD)
    \item The Linux kernel boot protocol documentation
          (\url{https://www.kernel.org/doc/html/v6.1/x86/boot.html})
    \item The SquashFS specification
          (\url{https://dr-emann.github.io/squashfs/})
    \item Binwalk (To see if anything major is missing)
\end{itemize}

% ==============================================================
\section{Image Layout}

Every byte in the image falls into one of the following regions. Exact
offsets are read from the headers themselves, no hardcoded assumptions
about layout are made.

\begin{lstlisting}[style=hexstyle, basicstyle=\ttfamily\footnotesize]
[0x00000000 			-> Live.UpgradeHeader]		Live Start Header + Metadata
[Live.UpgradeHeader 	-> Upgrade.MetaData.End] 	Upgrade Header + Metadata
[Upgrade.MetaData.End 	-> +0x200]            		Master Boot Record (512 bytes)
[MBR end 				-> kernelPayloadStart] 		Kernel Setup Sectors
[kernelPayloadStart 	-> kernelPayloadEnd]    	Kernel Payload (XZ vmlinux)
[kernelPayloadEnd 		-> initramfsStart]        	Unknown / Padding
[initramfsStart 		-> Live.FileSystemHeader]   Initramfs
[Live.FileSystemHeader 	-> Live.FileEnd]     		SquashFS Filesystem
\end{lstlisting}

% ==============================================================
\section{Region-by-Region Breakdown}

% --------------------------------------------------------------
\subsection{Live Start Header (\texttt{0x00000000})}

The image begins with a \textbf{Live Start Header}. This is a proprietary
header that acts as the outermost envelope for the entire firmware image.

\begin{center}
\begin{tabular}{llll}
\toprule
\textbf{Offset} & \textbf{Size} & \textbf{Field} & \textbf{Notes} \\
\midrule
\texttt{0x00} & 4  & Magic            & \\
\texttt{0x04} & 48 & Signature        & hash \\
\texttt{0x34} & 4  & Padding          & Always \texttt{0x00000000} \\
\texttt{0x38} & 8  & Anchor           & \texttt{DE AD BE EF DE AD BE EF} \\
\texttt{0x40} & 4  & UpgradeHeader    & Absolute offset to Upgrade Header \\
\texttt{0x44} & 4  & Unknown          & Purpose unknown \\
\texttt{0x48} & 4  & Flags            & \\
\texttt{0x4C} & 4  & FileSystemHeader & Absolute offset to SquashFS \\
\texttt{0x50} & 4  & FileEnd          & Absolute end of image \\
\texttt{0x54} & 56 & NullPad          & Zero padding \\
\bottomrule
\end{tabular}
\end{center}

After the fixed fields, a \textbf{metadata block} begins at
\texttt{0x8C} and extends to \texttt{UpgradeHeader}. This block uses a
TLV (Type-Length-Value) encoding with big-endian fields. It contains a
standard TLV section followed by a compact TLV section (1-byte tags,
2-byte lengths). The compact block ends with a 5-byte sentinel
(\texttt{0C 00 01 41 EB}) immediately before the next structural region.

The fields \texttt{UpgradeHeader}, \texttt{FileSystemHeader}, and
\texttt{FileEnd} fully determine where every subsequent region begins. Because this header starts at 0x0, these positions are also absolute offsets unlike every other offset in the file.

% --------------------------------------------------------------
\subsection{Upgrade Header}

Located at the absolute offset stored in \texttt{Live.UpgradeHeader}.

\begin{center}
\begin{tabular}{llll}
\toprule
\textbf{Offset} & \textbf{Size} & \textbf{Field} & \textbf{Notes} \\
\midrule
\texttt{0x00} & 4  & Magic            & \\
\texttt{0x04} & 48 & Signature        & \\
\texttt{0x34} & 4  & Padding          & \\
\texttt{0x38} & 8  & Anchor           & \texttt{DE AD BE EF DE AD BE EF} \\
\texttt{0x40} & 4  & MasterBootRecord & Relative offset to MBR \\
\texttt{0x44} & 4  & Unknown          & Purpose unknown \\
\texttt{0x48} & 4  & Flags            & \\
\texttt{0x4C} & 4  & Initramfs        & Relative offset to initramfs \\
\texttt{0x50} & 4  & FileSystemHeader & SquashFS offset (mirrors Live, relative to this header) \\
\texttt{0x54} & 56 & NullPad          & Zero padding \\
\bottomrule
\end{tabular}
\end{center}

Like the Live header, this is followed by its own metadata block
(TLV + compact TLV). The metadata block ends at
\texttt{Upgrade.MetaData.End}. These two headers are called Live and Upgrade because i suspect that the first one is the Live one actively being used and the second one is the one that is used when upgrading, though this is not confirmed.

% --------------------------------------------------------------
\subsection{Master Boot Record}

Begins immediately at \texttt{Upgrade.MetaData.End}. Standard 512-byte (0x200)
MBR layout:

\begin{center}
\begin{tabular}{lll}
\toprule
\textbf{Range} & \textbf{Size} & \textbf{Content} \\
\midrule
\texttt{0x000}--\texttt{0x1BD} & 446 & Bootstrap / boot code (x86) \\
\texttt{0x1BE}--\texttt{0x1FD} & 64  & Partition table (4 \(\times\) 16-byte entries) \\
\texttt{0x1FE}--\texttt{0x1FF} & 2   & Boot signature (\texttt{55 AA}) \\
\bottomrule
\end{tabular}
\end{center}

Each of the four partition entries contains:

\begin{itemize}
    \item \textbf{Status} (1 byte): \texttt{0x80} = bootable
    \item \textbf{CHS First / Last} (3 bytes each): legacy, mostly unused
    \item \textbf{LBA Start} and \textbf{Sector Count} (4 bytes each):
          the fields that actually describe the partition
    \item \textbf{Type} (1 byte): partition type identifier
\end{itemize}

% --------------------------------------------------------------
\subsection{Kernel Setup Sectors}

The kernel image begins immediately after the MBR
(\texttt{mbrStart + 0x200}). The Linux boot protocol defines a
\textbf{real-mode kernel header} at byte offset \texttt{0x1F1} within the
kernel image. Relevant fields:

\begin{center}
\begin{tabular}{lll}
\toprule
\textbf{Offset} & \textbf{Field} & \textbf{Notes} \\
\midrule
\texttt{0x1F1} & \texttt{setup\_sects} & Number of setup sectors \\
\texttt{0x200} & Jump              & \texttt{EB xx} short jump \\
\texttt{0x202} & Magic             & \texttt{HdrS} (\texttt{0x53726448}) \\
\texttt{0x206} & Version           & Boot protocol version \\
\texttt{0x248} & PayloadOffset     & Offset from setup end to compressed kernel \\
\texttt{0x24C} & PayloadLength     & Size of compressed kernel payload \\
\bottomrule
\end{tabular}
\end{center}

The setup region ends at \((\texttt{setup\_sects} + 1) \times 512\) bytes
from the kernel base. \texttt{PayloadOffset} is relative to that
boundary.

% --------------------------------------------------------------
\subsection{Kernel Payload}

\[
  \texttt{kernelPayloadStart} = \texttt{kernelStart} + \texttt{PayloadOffset}
\]
\[
  \texttt{kernelPayloadEnd} = \texttt{kernelPayloadStart} + \texttt{PayloadLength}
\]

This region contains the XZ-compressed \texttt{vmlinux}. It can be
decompressed independently if needed.

% --------------------------------------------------------------
\subsection{Unknown / Padding Region}

\[
  [\,\texttt{kernelPayloadEnd},\;
     \texttt{Live.UpgradeHeader} + \texttt{Upgrade.Initramfs}\,)
\]

This region lies between the end of the kernel payload and the start of
the initramfs. Its purpose was not fully determined, alignment padding
or an additional bootloader artifact is suspected.

% --------------------------------------------------------------
\subsection{Initramfs}

\[
  \texttt{initramfsStart} = \texttt{Live.UpgradeHeader} + \texttt{Upgrade.Initramfs}
\]
\[
  \texttt{initramfsEnd} = \texttt{Live.FileSystemHeader}
\]

The initramfs is a compressed CPIO archive. It is extracted as-is
(\texttt{initramfs.cpio}) in the output directory. The raw byte range is
fully determined by the header fields above.

% --------------------------------------------------------------
\subsection{SquashFS Filesystem}

Begins at the absolute offset stored in \texttt{Live.FileSystemHeader}.
This is the main read-only root filesystem. Parsing is described in full
in Section~\ref{sec:squashfs}.

% ==============================================================
\section{SquashFS Parsing}
\label{sec:squashfs}

The SquashFS specification at
\url{https://dr-emann.github.io/squashfs/} was followed directly.
The filesystem uses \textbf{XZ compression} (\texttt{CompressionID = 4}).

% --------------------------------------------------------------
\subsection{SquashFS Superblock (96 bytes, little-endian)}

\begin{center}
\begin{tabular}{llll}
\toprule
\textbf{Offset} & \textbf{Type} & \textbf{Field} & \textbf{Notes} \\
\midrule
\texttt{0x00} & \texttt{uint32} & Magic               & \texttt{0x73717368} \\
\texttt{0x04} & \texttt{uint32} & InodeCount          & \\
\texttt{0x08} & \texttt{uint32} & ModificationTime    & Unix timestamp \\
\texttt{0x0C} & \texttt{uint32} & BlockSize           & \\
\texttt{0x10} & \texttt{uint32} & FragmentEntryCount  & \\
\texttt{0x14} & \texttt{uint16} & CompressionID       & 4 = XZ \\
\texttt{0x16} & \texttt{uint16} & BlockLog            & \\
\texttt{0x18} & \texttt{uint16} & Flags               & \\
\texttt{0x1A} & \texttt{uint16} & IDCount             & \\
\texttt{0x1C} & \texttt{uint16} & VersionMajor        & \\
\texttt{0x1E} & \texttt{uint16} & VersionMinor        & \\
\texttt{0x20} & \texttt{uint64} & RootInodeRef        & Block + intra offset \\
\texttt{0x28} & \texttt{uint64} & BytesUsed           & \\
\texttt{0x30} & \texttt{uint64} & IDTableStart        & Absolute \\
\texttt{0x38} & \texttt{uint64} & XattrIDTableStart   & Absolute \\
\texttt{0x40} & \texttt{uint64} & InodeTableStart     & Relative to SquashFS base \\
\texttt{0x48} & \texttt{uint64} & DirectoryTableStart & Relative to SquashFS base \\
\texttt{0x50} & \texttt{uint64} & FragmentTableStart  & Absolute \\
\texttt{0x58} & \texttt{uint64} & ExportTableStart    & Absolute \\
\bottomrule
\end{tabular}
\end{center}

All table start offsets are used directly to locate each table.

% --------------------------------------------------------------
\subsection{Metadata Blocks}

All SquashFS metadata (inodes, directories, export table, ID table, xattr
table) is stored in \textbf{metadata blocks}. Each block is preceded by a
2-byte little-endian header:

\begin{itemize}
    \item \textbf{Bit 15}: \texttt{0} = XZ-compressed,
          \texttt{1} = uncompressed
    \item \textbf{Bits 14--0}: size of the following data in bytes
\end{itemize}

The \texttt{readSquashFSMetadataBlocks} function reads a contiguous range
of these blocks and returns the fully decompressed byte stream.

% --------------------------------------------------------------
\subsection{Inode Table}

Located between \texttt{InodeTableStart} and \texttt{DirectoryTableStart}
(both relative to the SquashFS base). Each inode begins with a common
16-byte header:

\begin{lstlisting}[style=gostyle]
type InodeHeader struct {
    InodeType    uint16
    Permissions  uint16
    UIDIdx       uint16
    GIDIdx       uint16
    ModifiedTime uint32
    InodeNumber  uint32
}
\end{lstlisting}

Followed by a type-specific body. Inode types parsed:

\begin{center}
\begin{tabular}{lll}
\toprule
\textbf{Type} & \textbf{Name} & \textbf{Notes} \\
\midrule
1  & Basic Directory    & Fixed 16-byte body \\
2  & Basic File         & Variable: \texttt{BlockSizes} array \\
3  & Basic Symlink      & Variable: target string \\
8  & Extended Directory & Variable: includes \texttt{DirIndex} array \\
9  & Extended File      & Variable: larger fields + \texttt{BlockSizes} \\
10 & Extended Symlink   & Includes \texttt{XattrIdx} \\
\bottomrule
\end{tabular}
\end{center}

For file inodes, the \texttt{BlockSizes} array length is computed as:

\[
  \texttt{blockCount} =
  \begin{cases}
    \left\lceil \dfrac{\texttt{FileSize}}{\texttt{BlockSize}} \right\rceil
      & \text{if } \texttt{FragmentBlockIndex} = \texttt{0xFFFFFFFF} \\[8pt]
    \left\lfloor \dfrac{\texttt{FileSize}}{\texttt{BlockSize}} \right\rfloor
      & \text{otherwise (tail stored in fragment)}
  \end{cases}
\]

% --------------------------------------------------------------
\subsection{Directory Table}

Located between \texttt{DirectoryTableStart} and
\texttt{FragmentTableStart}. Each directory consists of a 12-byte
\texttt{DirectoryHeader} followed by \(\texttt{Count} + 1\) entries.

\begin{lstlisting}[style=gostyle]
type DirectoryHeader struct {
    Count       uint32  // number of entries minus 1
    Start       uint32  // block offset into inode table
    InodeNumber uint32  // inode number of this directory
}

type DirectoryEntry struct {
    Offset      uint16
    InodeOffset int16   // signed delta from header InodeNumber
    Type        uint16
    NameSize    uint16
    Name        []uint8 // NameSize + 1 bytes
}
\end{lstlisting}

The absolute inode number for each entry is:
\[
  \texttt{inodeNumber} =
  \texttt{header.InodeNumber} + \texttt{entry.InodeOffset}
\]

% --------------------------------------------------------------
\subsection{Fragment Table}

For files whose final data does not fill a complete block, the remainder
is stored in a \textbf{fragment block}. The fragment table is an indirect
structure:

\begin{enumerate}
    \item \texttt{FragmentTableStart} points to a list of 8-byte absolute
          pointers (one per metadata block of fragment entries).
    \item Each pointer references a metadata block containing
          \texttt{FragmentBlockEntry} structs:
\end{enumerate}

\begin{lstlisting}[style=gostyle]
type FragmentBlockEntry struct {
    Start  uint64  // absolute offset of compressed fragment block
    Size   uint32  // bit 24: uncompressed flag; bits 23-0: size
    Unused uint32
}
\end{lstlisting}

To read a file's fragment tail: decompress the fragment block at
\texttt{Fragments[FragmentBlockIndex].Start}, then read
\texttt{remaining} bytes starting at \texttt{BlockOffset}.

% --------------------------------------------------------------
\subsection{Export Table}

Maps inode numbers to inode references (block offset + intra-block
offset). Used to locate the root inode by matching
\texttt{RootInodeRef}:

\begin{itemize}
    \item \textbf{Bits 63--16}: block offset from \texttt{InodeTableStart}
    \item \textbf{Bits 15--0}: byte offset within the decompressed block
\end{itemize}

% --------------------------------------------------------------
\subsection{ID Table}

Maps UID/GID index values (stored in inodes) to actual UID/GID integer
values. Uses the same indirect pointer-then-metadata-block structure as
the fragment table.

% --------------------------------------------------------------
\subsection{Xattr Table}

A three-level structure:

\begin{enumerate}
    \item \textbf{XattrIDTable} at \texttt{XattrIDTableStart}: contains
          \texttt{XattrTableStart} (absolute pointer to key/value data),
          \texttt{XattrIds} count, and a list of pointers to lookup
          metadata blocks.

    \item \textbf{Lookup table}: array of 16-byte
          \texttt{XattrLookupTable} structs, one per xattr set.
          Each entry contains an \texttt{XattrRef} (offset into the
          key/value data), a \texttt{Count} of key/value pairs, and a
          \texttt{Size}.

    \item \textbf{Key/value data} at \texttt{XattrTableStart}: compact
          TLV pairs. Each key entry is a 2-byte type, 2-byte name size,
          and name bytes. Each value entry is a 4-byte value size followed
          by value bytes. If bit 8 of the key type is set, the value is
          stored out-of-line.
\end{enumerate}

Each inode's \texttt{XattrIdx} field indexes into the lookup table to
find its associated xattr set.

% ==============================================================
\section{File Extraction}

Extraction walks every parsed inode. For each file inode:

\begin{enumerate}
    \item \textbf{Build the full path} by walking \texttt{ParentMap} from
          the inode up to the root inode, collecting names from
          \texttt{NameMap} at each step.

    \item \textbf{Write full data blocks} in order. For each entry in
          \texttt{BlockSizes}:
          \begin{itemize}
              \item Bit 24 set \(\Rightarrow\) uncompressed, read directly
              \item Bit 24 clear \(\Rightarrow\) XZ decompress
              \item Size = 0 \(\Rightarrow\) sparse block (zero-fill one
                    \texttt{BlockSize} of bytes)
          \end{itemize}

    \item \textbf{Write the fragment tail} if
          \texttt{FragmentBlockIndex} \(\neq\) \texttt{0xFFFFFFFF}:
          decompress the fragment block and append bytes
          \([\texttt{BlockOffset},\; \texttt{BlockOffset} +
          \texttt{remaining})\).
\end{enumerate}

Symlinks are created with \texttt{os.Symlink}. Directories are created
with \texttt{os.MkdirAll}. The \texttt{ParentMap} is built during
directory parsing by matching each directory inode's
\texttt{DirBlockStart} and \texttt{BlockOffset} against directory table
entries.

% ==============================================================
\section{What Remains Unexplained}

\begin{itemize}
    \item \textbf{The 5-byte compact TLV sentinel}
          (\texttt{0C 00 01 41 EB}): consistent across images, appears
          immediately before the next firmware region. Treated as a block
          terminator; semantic meaning is unknown.

    \item \textbf{\texttt{Unknown} in both headers}: a 4-byte big-endian
          value present in both Live and Upgrade headers. Possibly a size
          field.

    \item \textbf{Region between kernel payload and initramfs}:
          It is the Linux Kernel Early Boot Stub and Decompressor (specifically for x86-64).

    \item \textbf{Metadata block \texttt{BlockType} field}: parsed but not
          interpreted, likely a version or type discriminator for the
          proprietary header format.
\end{itemize}

\end{document}